\section{Dirichlet's and Neumann's problems on a plane that crosses the sources}
\label{sec:generalized_dirichlet_neumann}


It is important noting that the integral equations \ref{eq:U-Dirichlet} and 
\ref{eq:U-Neumann} are valid for the case in which the planar layer of monopoles
is defined on a plane $z^{\dagger}_{0}$ that does not cross the actual 3D sources. 
What follows aims at generalizing these integral equations for the case in which the 
monopole layer is defined on a plane that crosses the 3D sources.
To do this, let us first consider the following integral equation:
\begin{equation}
	U(\mathbf{r}^{\dagger}) = 
	\int\limits_{-\infty}^{\infty} \int\limits_{-\infty}^{\infty}
	\tau(\mathbf{r}'') \:
	\partial_{z} \Psi(\mathbf{r}^{\dagger}, \mathbf{r}'') 
	\:\: dx''dy'' \: ,
	\label{eq:U-tau}
\end{equation}
where $\partial_{z} \Psi(\mathbf{r}^{\dagger}, \mathbf{r}'')$ is the vertical
derivative of the same scaled inverse distance
defined by equation \ref{eq:Psi}, but in this case computed for a point
$\mathbf{r}^{\dagger} = (x^{\dagger}, y^{\dagger}, z^{\dagger}_{0})^{\top}$ on the
horizontal plane $z^{\dagger}_{0}$ above the 3D sources and a point 
$\mathbf{r}'' = (x'', y'', z''_{0})^{\top}$ 
on a horizontal plane that is located at $z''_{0} = z^{\dagger}_{0} + \Delta z$, 
with $\Delta z > 0$ , and may cross the 3D sources.
By applying the Fourier-Convolution Theorem \cite[][ p. 268]{blakely1996},
we obtain the following closed-form solution for $\tau(\mathbf{r}'')$ in Fourier domain:
\begin{equation}
	\begin{split}
		\mathcal{F} \left[ \tau(\mathbf{r}'') \right] 
		&= \mathcal{F} \left[ U(\mathbf{r}^{\dagger}) \right] {\Big /}
		\mathcal{F} \left[ \partial_{z} \Psi(\mathbf{r}^{\dagger}, \mathbf{r}'') \right] \\
		&= \mathcal{F} \left[ U(\mathbf{r}^{\dagger}) \right] 
		\, e^{\vert k \vert \, \Delta z}
	\end{split} \quad , \quad \Delta z > 0 \: , \quad \vert k \vert \ne 0 \: ,
	\label{eq:Fourier-U-tau}
\end{equation}
where $\mathcal{F}\left[ \cdot \right]$ denotes the 2D Fourier Transform (2DFT)
with respect to the horizontal coordinates $x$ and $y$.
The 2DFT's are two-dimensional functions of the wavenumbers $k_{x}$ and $k_{y}$ 
along the $x$ and $y$ directions, respectively.
The wavenumbers have units of inverse distance. In equation \ref{eq:Fourier-U-tau},
$k = \sqrt{k_{x}^{2} + k_{y}^{2}}$, the 2DFT of 
$\partial_{z} \Psi(\mathbf{r}^{\dagger}, \mathbf{r}'')$ is computed according to
\citet[][ p. 275]{blakely1996} and the dependence of the 2DFT's on $k_{x}$ and $k_{y}$
is conveniently omitted to simplify the notation.
From a geophysical point of view, $\tau(\mathbf{r}'')$ represents a  
potential field obtained by computing the downward continuation of 
$U(\mathbf{r}^{\dagger})$ to points that may be located inside the sources. 
This potential field is different from the $U(\mathbf{r}'')$ produced by the true 3D 
sources on the plane $z''_{0}$ for the case in which this plane crosses the sources.
Note that $\tau(\mathbf{r}'')$ is equal to $U(\mathbf{r}'')$ for the case in which the
horizontal plane $z''_{0}$ is located above the sources.

Now, let us substitute the integral equation \ref{eq:U-tau} into the solution of the
Dirichlet's problem (equation \ref{eq:U-Dirichlet}) and change the integration order:
\begin{equation}
	\begin{split}
		U(\mathbf{r}) 
		&= \int\limits_{-\infty}^{\infty} \int\limits_{-\infty}^{\infty}
		\left[ 
		\int\limits_{-\infty}^{\infty} \int\limits_{-\infty}^{\infty}
		\tau(\mathbf{r}'') \:\: 
		\partial_{z} \Psi(\mathbf{r}^{\dagger}, \mathbf{r}'') 
		\: dx''dy'' \right]
		\partial_{z} \Psi(\mathbf{r}, \mathbf{r}^{\dagger}) 
		\:\: dx^{\dagger}dy^{\dagger} \\
		&= 
		\int\limits_{-\infty}^{\infty} \int\limits_{-\infty}^{\infty}
		\tau(\mathbf{r}'') \:\:
		\left[ 
		\int\limits_{-\infty}^{\infty} \int\limits_{-\infty}^{\infty}
		\partial_{z} \Psi(\mathbf{r}^{\dagger}, \mathbf{r}'')
		\: \partial_{z} \Psi(\mathbf{r}, \mathbf{r}^{\dagger}) 
		\: dx^{\dagger}dy^{\dagger} \right]
		\:\: dx''dy'' \quad .
	\end{split}
	\label{eq:intermediary-equation}
\end{equation}
Finally, by noting that the term within brackets in the second line of equation
\ref{eq:intermediary-equation} represents the solution of Dirichlet's problem
for the function $\partial_{z} \Psi(\mathbf{r}, \mathbf{r}'')$ on the plane
$z^{\dagger}_{0}$, we obtain:
\begin{equation}
	U(\mathbf{r}) =  
	\int\limits_{-\infty}^{\infty} \int\limits_{-\infty}^{\infty}
	\tau(\mathbf{r}'') \:\: 
	\partial_{z} \Psi(\mathbf{r}, \mathbf{r}'') 
	\:\: dx''dy'' \: , \quad z < z''_{0} \: .
	\label{eq:U-Dirichlet-cross}
\end{equation}
This equation shows that the same continuous function $\tau(\mathbf{r}'')$
that reproduces the potential $U(\mathbf{r}^{\dagger})$ on the horizontal plane 
$z^{\dagger}_{0}$ above the sources also reproduces the potential $U(\mathbf{r})$ 
at all observation points 
$\mathbf{r} = (x, y, z)^{\top}$ located above $z^{\dagger}_{0}$, with $z < z^{\dagger}_{0}$.
Equation \ref{eq:U-Dirichlet-cross} generalizes the solution of Dirichlet's problem
for the case in which the horizontal plane $z''_{0}$ crosses the true 3D sources.

Consider now the following integral equation:
\begin{equation}
	U(\mathbf{r}^{\dagger}) = 
	\int\limits_{-\infty}^{\infty} \int\limits_{-\infty}^{\infty}
	\upsilon(\mathbf{r}'') \:
	\Psi(\mathbf{r}^{\dagger}, \mathbf{r}'') 
	\:\: dx''dy'' \: .
	\label{eq:U-upsilon}
\end{equation}
Then, by applying the Fourier-Convolution Theorem \cite[][ p. 268]{blakely1996},
we obtain the following closed-form solution for 
$\upsilon(\mathbf{r}'')$ in Fourier domain:
\begin{equation}
	\begin{split}
		\mathcal{F} \left[ \upsilon(\mathbf{r}'') \right] 
		&= \mathcal{F} \left[ U(\mathbf{r}^{\dagger}) \right] {\Big /}
		\mathcal{F} \left[ \Psi(\mathbf{r}^{\dagger}, \mathbf{r}'') \right] \\
		&= \mathcal{F} \left[ U(\mathbf{r}^{\dagger}) \right] 
		\vert k \vert \, e^{\vert k \vert \, \Delta z} \\
		&= \mathcal{F} \left[ \partial_{z} U(\mathbf{r}^{\dagger}) \right] 
		\, e^{\vert k \vert \, \Delta z} \\
	\end{split} \quad , \quad \Delta z > 0 \: , \quad \vert k \vert \ne 0 \: ,
\label{eq:Fourier-U-upsilon}
\end{equation}
From a geophysical point of view, $\upsilon(\mathbf{r}'')$ represents a  
potential field obtained by computing the downward continuation of 
$\partial_{z} U(\mathbf{r}^{\dagger})$ to points that may be located inside the sources. 
This potential field is different from the $\partial_{z} U(\mathbf{r}'')$ produced by 
the true 3D sources on the plane $z''_{0}$ for the case in which this plane crosses 
the sources.
Note that $\upsilon(\mathbf{r}'')$ is equal to $\partial_{z} U(\mathbf{r}'')$ for the 
case in which the horizontal plane $z''_{0}$ is located above the sources.

Now, let us substitute the vertical derivative of the integral equation \ref{eq:U-upsilon}
into the solution of the Neumann's problem (equation \ref{eq:U-Neumann}) and change the
integration order:
\begin{equation}
	\begin{split}
		U(\mathbf{r}) 
		&= \int\limits_{-\infty}^{\infty} \int\limits_{-\infty}^{\infty}
		\left[ 
		\int\limits_{-\infty}^{\infty} \int\limits_{-\infty}^{\infty}
		\upsilon(\mathbf{r}'') \:\: 
		\partial_{z} \Psi(\mathbf{r}^{\dagger}, \mathbf{r}'') 
		\: dx''dy'' \right]
		\Psi(\mathbf{r}, \mathbf{r}^{\dagger}) 
		\:\: dx^{\dagger}dy^{\dagger} \\
		&= 
		\int\limits_{-\infty}^{\infty} \int\limits_{-\infty}^{\infty}
		\upsilon(\mathbf{r}'') \:\:
		\left[ 
		\int\limits_{-\infty}^{\infty} \int\limits_{-\infty}^{\infty}
		\partial_{z} \Psi(\mathbf{r}^{\dagger}, \mathbf{r}'')
		\: \Psi(\mathbf{r}, \mathbf{r}^{\dagger}) 
		\: dx^{\dagger}dy^{\dagger} \right]
		\:\: dx''dy'' \quad .
	\end{split}
\label{eq:intermediary-equation-2}
\end{equation}
Finally, by noting that the term within brackets in the second line of equation
\ref{eq:intermediary-equation-2} represents the solution of Neumann's problem
for the function $\Psi(\mathbf{r}, \mathbf{r}'')$ on the plane
$z^{\dagger}_{0}$, we obtain:
\begin{equation}
U(\mathbf{r}) =  
\int\limits_{-\infty}^{\infty} \int\limits_{-\infty}^{\infty}
\upsilon(\mathbf{r}'') \:\: 
\Psi(\mathbf{r}, \mathbf{r}'') 
\:\: dx''dy'' \: , \quad z < z''_{0} \: .
\label{eq:U-Neumann-cross}
\end{equation}
Equation \ref{eq:U-Neumann-cross} shows that the same continuous function
$\upsilon(\mathbf{r}'')$ that reproduces the potential $U(\mathbf{r}^{\dagger})$ on 
the horizontal plane $z^{\dagger}_{0}$ above the sources 
also reproduces the potential $U(\mathbf{r})$ at all
observation points $\mathbf{r} = (x, y, z)^{\top}$ located above $z^{\dagger}_{0}$, 
with $z < z^{\dagger}_{0}$.
Equation \ref{eq:U-Neumann-cross} generalizes the solution of Neumann's problem
for the case in which the horizontal plane $z''_{0}$ crosses the true 3D sources.